\documentclass[varwidth=false]{standalone}

% ======== Reproducimos estilo copernicus.cls ========

\usepackage[T1]{fontenc}
\usepackage{mathptmx}                 % Serif con soporte matemático (Times)
\usepackage{times}
\renewcommand{\familydefault}{\rmdefault} % Helvetica como texto por defecto, como en copernicus.cls
\usepackage{courier}                 % Monoespaciada
\usepackage{booktabs}                % Tablas con líneas finas
\usepackage{siunitx}                 % Para alineación numérica
\renewcommand{\baselinestretch}{1.2}

% Ajuste fino de grosor de líneas en tablas
\setlength{\arrayrulewidth}{0.4pt}
\setlength{\tabcolsep}{6pt}

\begin{document}

\begin{tabular}{|l|l|l|l|r|}

\hline
\textbf{Variable} & \textbf{Magnitud} & \textbf{Fuente} & \textbf{Estación}	& \textbf{Unidad}\\
\hline
$P_\mathrm{MAQ}$ & Precipitación acumulada & DMC & Maquehue (Ad.) & mm\\
\hline
$Q_\mathrm{BCC}$ & Caudal instantáneo & DGA & Río Blanco en Curacautín & m³/s\\
\hline
$Q_\mathrm{CJN}$ & Caudal instantáneo & DGA & Río Cautín en Cajón & m³/s\\
\hline
$Q_\mathrm{CRR}$ & Caudal instantáneo & DGA & Río Cautín en Rari-Ruca & m³/s\\
\hline


\end{tabular}

\end{document}